\setcounter{section}{8}
\section{Recursos y referencias de contraste}
Se inventariaron los archivos del repositorio y se comprobaron los tres ZIP existentes medianteCRC. ElZIP de Capachica se abrió y sus seis DXF se analizaron; aporta organización planta/perfil/unifilar/detalles, no prueba las cargas de Limbani. ElZIPLP-R01 se conserva íntegro; incluye archivos de bloqueoDW L/DWL2 que no se incorporan al paquete nuevo. La bibliotecaDGE es referencia; no constituye adaptación validada aCAC.

Los términos de Electro Puno consultados utilizan once volúmenes y anexos. Las ocho carpetas aquí responden a organización académica, no a una exigencia universal de ocho capítulos. Se agrupan cálculos eléctricos/mecánicos/protección en 03, armados en 04 y especificaciones en 05, preservando la correspondencia de contenido.

SEACE: se descargaron dos contratos oficiales comparables, Ubinas33 kV yCarhuaz13.2 kV, y la ficha de poste 13/400. Los contratos no incluyen los planos completos del expediente. Para modelos gráficos se descargaron del SINIA el planos de componentesLircay, mapaTablachaca y planillaChoclococha. El buscadorSEACE redirige a una aplicación dinámica; los documentos directos pudieron recuperarse por descargaHTTP aunque la herramienta web devolvía 403.

\begingroup\small\setlength{\LTpre}{4pt}\setlength{\LTpost}{6pt}
\begin{longtable}{*{2}{L{\dimexpr\linewidth/2-2\tabcolsep\relax}}}
\toprule
\textbf{ID y entidad} & \textbf{Documento local / límite} \\ \midrule
\endfirsthead
\toprule
\textbf{ID y entidad} & \textbf{Documento local / límite} \\ \midrule
\endhead
R02-F01 MINEM & rd018\_\allowbreak{}2003.pdf \ensuremath{\cdot} Bases primarias; texto aprobado 2003. No sustituye consolidación normativa completa. \\
R02-F02 MINEM & cne\_\allowbreak{}suministro\_\allowbreak{}2011.pdf \ensuremath{\cdot} Cargas, grados, distancias, factores. Revisar modificaciones aplicables al cierre. \\
R02-F03 SEACE/PERU COMPRAS & poste\_\allowbreak{}seace\_\allowbreak{}13\_\allowbreak{}400.pdf \ensuremath{\cdot} Ficha de bien 13/400/2/180/375; rotulado daN. No certificado de poste 600 daN ni ensayo axial. \\
R02-F04 Electro Puno & tdr\_\allowbreak{}electropuno\_\allowbreak{}lp.pdf \ensuremath{\cdot} Referencia de estructura documental Tototani; normativa histórica requiere cotejo actual. \\
R02-F05 SEACE/Electrosur & seace\_\allowbreak{}ubinas\_\allowbreak{}contrato\_\allowbreak{}023\_\allowbreak{}2021.pdf \ensuremath{\cdot} LP007-2020-ES. Línea33 kV. Contrato y resumen de presupuesto; no planos de ingeniería. \\
R02-F06 SEACE/Hidrandina & seace\_\allowbreak{}carhuaz\_\allowbreak{}contrato\_\allowbreak{}091\_\allowbreak{}2024.pdf \ensuremath{\cdot} AS-SM44-2023-HDNA-3. Consultoría13.2 kV; no expediente técnico adjunto descargado. \\
R02-F07 SINIA/Municipalidad Huando & dia\_\allowbreak{}tinyacclla\_\allowbreak{}pachachaca.pdf \ensuremath{\cdot} Referencia de componentes. Se observaron diferencias 13.2/22.9 kV y presupuesto: requiere evaluación crítica. \\
R02-F08 SINIA & dia-hospital\_\allowbreak{}lircay.pdf \ensuremath{\cdot} DIA de sistema 22.9 kV; se revisaron planos decomponentes/trazo/ubicación en anexos186--195. No trasladar 1250 kVA ni ajustes. \\
R02-F09 SINIA & plano\_\allowbreak{}trazo\_\allowbreak{}tablachaca.pdf \ensuremath{\cdot} Plano de ejes con coordenadas y cuadro de tramos. Referencia de presentación, no ruta de Limbani. \\
R02-F10 SINIA & dia\_\allowbreak{}choclococha.pdf \ensuremath{\cdot} Planilla de estructuras con poste/armado/vano/coordenadas. No transferir clima ni equipos. \\
\bottomrule\end{longtable}\endgroup

Cada fuente conserva URL, fecha, SHA-256, tamaño y uso permitido en 08-recursos/fuentes.json. Las referencias históricas pueden tener inconsistencias: se observaron diferencias de tensión/totales enTinyacclla y porcentajes del resumen deUbinas; no se trasladan sus cifras como reglas o precios actuales.

